Tras la validación del prototipo funcional \textbf{ELCO-Dealer}, se han identificado diversas líneas de evolución tecnológica que permitirían transformar este desarrollo académico en un producto de electrónica de consumo líder en el mercado masivo.

\begin{itemize}
    \item \textbf{Evolución de la Interfaz de Usuario (UI):} Se plantea la sustitución del actual LCD Keypad Shield por una \textbf{pantalla táctil capacitiva}. Esto permitiría una navegación más fluida y la implementación de una \textbf{interfaz inteligente potenciada por IA}, capaz de sugerir modos de reparto y llevar estadísticas de juego en tiempo real.

    
    \item \textbf{Conectividad IoT y Aplicación Móvil:} El sistema podría integrar un módulo \textbf{Bluetooth o Wi-Fi} para sincronizarse con una aplicación móvil dedicada. Desde esta app, los usuarios podrían monitorizar el número de cartas restantes, configurar modos de reparto personalizados para juegos específicos y gestionar puntuaciones automáticamente mediante el uso de visión artificial.
    
    \item \textbf{Optimización de Hardware e Industrialización:} Siguiendo la estrategia comercial para 10.000 unidades, se proyecta la transición de módulos comerciales a una \textbf{PCB propia (System-on-Board)}. Esto permitiría eliminar el exceso de cableado, integrar los drivers de los motores directamente como chips y reducir el consumo energético del conjunto, optimizando el espacio y los costes de fabricación.
\end{itemize}